\documentclass[11pt,titlepage]{article}

\usepackage[french,american]{babel}
\usepackage[utf8]{inputenc}
\usepackage[T1]{fontenc}
\usepackage{lmodern}
\usepackage{amsmath,amsfonts,amssymb}
\usepackage{graphicx}
\usepackage[a4paper]{geometry}           % See geometry.pdf to learn the layout
%\geometry{landscape}                % Activate for for rotated page geometry
%\usepackage[parfill]{parskip}    % Activate to begin paragraphs with an empty line rather than an indent
\usepackage{graphicx}
\usepackage{amssymb}
\usepackage{epstopdf}
\usepackage{color}
\usepackage[tt]{titlepic}
\usepackage{fancyhdr}
\usepackage{secdot}
\usepackage{booktabs}
\usepackage{rotating}
\usepackage{multirow}
\usepackage{tabularx}
\usepackage{ragged2e}  % for '\RaggedRight' macro (allows hyphenation)
\newcolumntype{Y}{>{\RaggedRight\arraybackslash}X} 
\usepackage[style=ieee, backend=bibtex]{biblatex}

\topmargin=-0.45in      %
%\evensidemargin=0in     %
\oddsidemargin=-0.1in      %
\textwidth=6.8in        %
\textheight=9.0in       %
%\headsep=0.25in         %

%======================= Titlepage =========================================================
\DeclareGraphicsRule{.tif}{png}{.png}{`convert #1 `dirname #1`/`basename #1 .tif`.png}
\titlepic{\includegraphics[width=17cm]{figures/research_plan_pic.pdf}}

\title{\textbf{Annual progress report \textcolor{blue}{2014}}} %Change Year

\author{\large{\textbf{Damar Wicaksono} }\\\\ 
\large{\textbf{Prof. Andreas Pautz}} \\\
\large{\textbf{Omar Zerkak}} \\\\\\\
\Large{\textbf{Bayesian Uncertainty Quantification}} \\\\
\Large\textbf{of Physical Models} \\\\
\Large\textbf{in Thermal-Hydraulics System Code}}

\date{30.11.2014}				%activate to remove date
%========================================================================================

\bibliography{report_2014.bib}

\begin{document}
%========================  Top header ======================================================
\pagestyle{fancy} \pagenumbering{arabic} \setcounter{page}{1}
\addtolength{\headheight}{\baselineskip}
\newcommand{\ffont}{\fontsize{8}{8}\selectfont}
\lhead{ECOLE DOCTORALE\\ \textit{DOCTORAL SCHOOL}}
\chead{PROGRAMME DOCTORAL EN PHYSIQUE\\ \textit{DOCTORAL PROGRAM IN PHYSICS}}
\rhead{\bfseries\ffont  \includegraphics[width=55pt]{figures/EPFL_LOG.pdf}}
\renewcommand{\headrulewidth}{0.4pt}
\maketitle
%=========================================================================================
\textcolor{white}.
\section{Objectives of research}

The objective of the research is to quantify the uncertainty of 
physical model parameters implemented in a thermal-hydraulics system code.
The physical models concerned are the ones describing 
the phasic interaction (heat and momentum exchange) in a complex 
multiphase flow during reactor transient especially for the
\textit{reflood} phase after a loss-of-coolant-accident. 
These models parameterized either by a physical or tuning 
parameters.

Conforming with the practice of statistical uncertainty propagation widely 
adopted in the field of nuclear engineering, probability theory is used to 
quantify the uncertainties related to these parameters with the results in 
a form of density function and/or its approximation. 
The derivation of this density is posed as an inverse statistical problem 
following a bayesian approach as the parameters themselves are not directly 
observable.
To this end, a methodology of model uncertainty quantification will be 
developed combining on probabilistic modeling and the available relevant 
experimental data taken from various separate effect test facilities.

\section{Work achieved in the past year (state of research)}

\subsection{Assessment of TRACE on the basis of FEBA and PERICLES Test Facilities}

\subsection{PSI Contribution to OECD/NEA PREMIUM Benchmark}

\subsection{Application of Morris Screening Method}

\subsection{Application of Functional Data Analysis}

\section{Current state of work}

The current state of research along with the achievements during the year are 
summarized in the table below.
                   analysis and
                               novel data analysis methodology (NUTHOS-14)
                          1 abstract for the NURETH-16 submitted 
                                             
\begin{table}[!h]
\caption{Current state of research in relation to the previously submitted 
         (1\textsuperscript{st} year) work plan
         } 
\begin{tabularx}{\textwidth}{@{} c Y Y Y @{}} % use 'Y' for first column
% Header ----------------------------------------------------------------------
    \toprule
    \textbf{Phase} 
 & \textbf{Task Description} 
 & \textbf{Planned Outcome} 
 & \textbf{Current State of Work} \\ \midrule
% Table Elements --------------------------------------------------------------
   \multirow{5}{*}{\textbf{1}}                                %Phase
 & Comprehensive reviews of post-CHF flow closure             %Task Description
   laws in TRACE and externalization of important 
   model parameters  
 & 1 Technical Report (PSI)                                   %Planned Outcome
 & 1. Important reflood model parameters externalized         %Current State
   2. PSI contribution to OECD/NEA PREMIUM Benchmark finalized  \\ \midrule
%
   \multirow{8}{*}{\textbf{2}}                                %Phase
 & Global sensitivity analysis of important                   %Task Description
   based on FEBA test facility 
 & 1 Technical Report (PSI) \& 1 Journal Article              %Planned Outcome
 & 1. A paper on Morris screening method for NUTHOS-14        %Current State
   2. Another paper on an application of functional data 
      analysis for NUTHOS-14
   3. One abstract on global sensitivity analysis 
      method submitted for NURETH-16      \\ \midrule
%
   \multirow{5}{*}{\textbf{3}}                                %Phase
 & Calibration of TRACE reflood model based of FEBA test      %Task Description 
   facility; Definition of error and specification of 
   probabilistic error model;
 & 1 Technical Report (PSI) \& 1 Journal Article              %Planned Outcome
 & $-$ \\ \midrule                                            %Current State
%
   \multirow{6}{*}{\textbf{4}}
 & Calibration of TRACE reflood model based on another 
   reflood test facility 
 & 1 Conference Paper 
 & Possible candidates of reflood facility with adequate 
 specification and 
   experimental data are made available (\textit{e.g.,} 
   SEFLEX, NEPTUN, ACHILLES) \\ \midrule
%
   \multirow{4}{*}{\textbf{5}}
 & Consolidation of the calibration results based on 2 facilities and 
   validation based on another reflood test facility 
 & 1 Journal Article 
 & $-$  \\
   \textbf{6} 
 & Thesis write-up  
 & Thesis          
 & $-$ \\
\bottomrule
\end{tabularx}
\label{table:nonlin}
\end{table}


\section{Calendar of upcoming work}


\pagebreak

\section{Scientific publications, conference contributions, and presentations}

\nocite{Wicaksono2014a}
\nocite{Wicaksono2014b}
\nocite{Wicaksono2014c}
\nocite{Wicaksono2014d}
\nocite{Wicaksono2014e}
\printbibliography[heading=none]

\section{Other activities and remarks}

\subsection{NES PhD Student Day}

\subsection{Nuclear Computation Lab (ETH-531)}

A teaching activity during the year 2014 was carried out for Nuclear Computation Lab 
course given at PSI.
The course was part of the block courses offered to the EPFL/ETHZ Nuclear 
Engineering Master Program
Single (full) day was dedicated for lecture session followed by computer 
simulation session for the t
12 students participated in the course with 

In preparation for the course, the available lecture slides was fully revised 
for a better structure and with additional examples.
The simulation session was also slightly extended with more freedom given to 
the students.

Overall impression of the students


 

\end{document}
\newpage
\textcolor{white}.\\\\
\noindent \textbf{Comments by the thesis advisor(s)}\\\\
...................................................................................................................................................................\\\\
...................................................................................................................................................................\\\\
...................................................................................................................................................................\\\\
...................................................................................................................................................................\\\\
...................................................................................................................................................................\\\\
...................................................................................................................................................................\\\\
...................................................................................................................................................................\\\\
...................................................................................................................................................................\\\\
...................................................................................................................................................................\\\\


\section*{\underline{Signatures}\\}
\noindent \textbf{Thesis advisor}\hspace{6.25cm}..............................................................................\vspace{0.5cm}\\

\noindent \textbf{Thesis co-advisor}\hspace{5.7cm}..............................................................................\\
\noindent  (if officially nominated)\\\\

\noindent \textbf{With his/her signature, the candidate confirms that he/she took note of the above comments.}\\\\
\textbf{Candidate}\hspace{7cm}..............................................................................\vspace{0.5cm}\\

\noindent \textbf{Date}\hspace{8.05cm}..............................................................................\\

\vspace{0.8cm}
\begin{center}
\textbf{Up to 5 pages in total}
\end{center}
\begin{center}
\textbf{To be returned to:} \textcolor{green}{EPFL-EDPY / Station 3 / 1015 Lausanne / Suisse}
\end{center}
\end{document}  